\editionPage[\pdfbookmark[0]{Corrigés des bilans}{bm-bilan-A-cor-auto-parent}]{bilan-A-cor-auto}{Bilan A · Corrigé}{Automatismes expliqués}{\hyperref[bilan-A-auto]{Énoncé p.\,\pageref*{bilan-A-auto}}}
\begin{blocEdition}{bilan-A-cor-auto}
\begin{questions}
\item $\res{8\,307}$ : chaque valeur rejoint son rang.
\item $457+20-1=\res{476}$.
\item $\res{244}$ ; contrôle : $244+458=702$.
\item $\res9$, car $9\times4=36$.
\item $7\times6=42$ puis un zéro : $\res{420}$.
\item $\res{3\,400<3\,456<3\,500}$.
\item $\res{\frac58}$ : les huitièmes s’ajoutent.
\item $\res{2{,}35}$.
\item $4{,}50>4{,}45$ : $\res{4{,}5>4{,}45}$.
\item $\res{36}$ : chaque chiffre monte d’un rang.
\item $\square=\res4$.
\item $\res7$, car $6{,}78$ est plus proche de 7.\end{questions}
\end{blocEdition}
\editionPage{bilan-A-cor-1}{Bilan A · Corrigé 1}{La fête de l’école}{\hyperref[bilan-A-ex-1]{Énoncé p.\,\pageref*{bilan-A-ex-1}}}
\begin{blocEdition}{bilan-A-cor-1}
\editionEtape{Recette}$875+468=\res{1\,343}$ euros.\editionEtape{Reste}$1\,343-390=\res{953}$ euros.\editionEtape{Partage}$953=7\times136+1$ : chaque classe reçoit $\res{136}$ euros et il reste 1 euro. Contrôle : $7\times136+1=953$.
\end{blocEdition}
\editionPage{bilan-A-cor-2}{Bilan A · Corrigé 2}{Les cartons de la médiathèque}{\hyperref[bilan-A-ex-2]{Énoncé p.\,\pageref*{bilan-A-ex-2}}}
\begin{blocEdition}{bilan-A-cor-2}
\editionEtape{Division}$1\,350=8\times168+6$ : $\res{168}$ cartons complets.\editionEtape{Le reste}Il reste $\res6$ livres.\editionEtape{Interpréter}Pour tout transporter, un carton entamé s’ajoute : $\res{169}$ cartons. Le même reste se lit selon la question.
\end{blocEdition}
\editionPage{bilan-A-cor-3}{Bilan A · Corrigé 3}{Le gâteau du dimanche}{\hyperref[bilan-A-ex-3]{Énoncé p.\,\pageref*{bilan-A-ex-3}}}
\begin{blocEdition}{bilan-A-cor-3}
\editionEtape{Après le midi}$\frac88-\frac58=\res{\frac38}$ du gâteau.\editionEtape{Sur la journée}$5+3=8$ parts : $\res{\frac88=1}$, tout le gâteau.\editionEtape{Le second gâteau}Chaque enfant reçoit $\res{\frac14}$ de gâteau, soit 2 parts sur 8 : $\frac28=\frac14$.
\end{blocEdition}
\editionPage{bilan-A-cor-4}{Bilan A · Corrigé 4}{Les courses du mercredi}{\hyperref[bilan-A-ex-4]{Énoncé p.\,\pageref*{bilan-A-ex-4}}}
\begin{blocEdition}{bilan-A-cor-4}
\editionEtape{Dépense}$8{,}70+4{,}55=\res{13{,}25}$ euros.\editionEtape{Monnaie}$20-13{,}25=\res{6{,}75}$ euros.\editionEtape{Comparer}Trois carnets coûtent $2{,}20\times3=6{,}60$ euros et $6{,}60<6{,}75$ : $\res{\text{oui}}$, il resterait 15 centimes.
\end{blocEdition}
\editionPage{bilan-B-cor-auto}{Bilan B · Corrigé}{Automatismes expliqués}{\hyperref[bilan-B-auto]{Énoncé p.\,\pageref*{bilan-B-auto}}}
\begin{blocEdition}{bilan-B-cor-auto}
\begin{questions}
\item $\res8$.
\item $680:2=\res{340}$.
\item $801-600+1=\res{202}$.
\item $75=8\times9+3$ : quotient $\res9$, reste $\res3$.
\item $\res{1+\frac34}$.
\item $\frac{75}{100}=\res{\frac34}$.
\item $\res{5{,}08<5{,}8}$.
\item $\res{6{,}55}$ par compléments.
\item $120+15=\res{135}$ min.
\item $4\times9=\res{36}$ cm.
\item $3\,000>2\,900$ : $\res{3\ \mathrm{kg}>2\,900\ \mathrm g}$.
\item Deux fois plus : $\res{12}$ euros.\end{questions}
\end{blocEdition}
\editionPage{bilan-B-cor-1}{Bilan B · Corrigé 1}{Le train de Lucie}{\hyperref[bilan-B-ex-1]{Énoncé p.\,\pageref*{bilan-B-ex-1}}}
\begin{blocEdition}{bilan-B-cor-1}
\editionEtape{L’aller}De 9 h 47 à 10 h : 13 min ; puis 2 h ; puis 15 min : $\res{2\ \text h\ 28}$.\editionEtape{Le retour}Reculer de 2 h 40 depuis 19 h 05 : $19\ \text h\ 05\to17\ \text h\ 05\to16\ \text h\ 25$ : départ à $\res{16\ \text h\ 25}$.\editionEtape{Comparer}$2$ h $40=160$ min et $2$ h $28=148$ min : le retour dépasse l’aller de $\res{12}$ min.
\end{blocEdition}
\editionPage{bilan-B-cor-2}{Bilan B · Corrigé 2}{Le potager de l’école}{\hyperref[bilan-B-ex-2]{Énoncé p.\,\pageref*{bilan-B-ex-2}}}
\begin{blocEdition}{bilan-B-cor-2}
\editionEtape{Périmètre}$2\times(7{,}5+4)=2\times11{,}5=\res{23}$ m.\editionEtape{Prix}Un mètre coûte $2{,}50$ euros : 23 m coûtent $2{,}50\times23=\res{57{,}50}$ euros, car $2{,}5\times23=57{,}5$.\editionEtape{Reste}$60-57{,}50=\res{2{,}50}$ euros. Contrôle : le prix est proportionnel à la longueur achetée.
\end{blocEdition}
\editionPage{bilan-B-cor-3}{Bilan B · Corrigé 3}{L’enquête des déplacements}{\hyperref[bilan-B-ex-3]{Énoncé p.\,\pageref*{bilan-B-ex-3}}}
\begin{blocEdition}{bilan-B-cor-3}
\editionEtape{Colonne bus}$9+13=\res{22}$ élèves.\editionEtape{Comparer}12 marcheurs en CM1 contre 10 en CM2 : la marche domine en $\res{\text{CM1}}$.\editionEtape{Contrôle global}$12+9+5=26$ en CM1 et $10+13+6=29$ en CM2 : $26+29=\res{55}$ élèves, le compte est bon.
\end{blocEdition}
\editionPage{bilan-B-cor-4}{Bilan B · Corrigé 4}{La kermesse du hasard}{\hyperref[bilan-B-ex-4]{Énoncé p.\,\pageref*{bilan-B-ex-4}}}
\begin{blocEdition}{bilan-B-cor-4}
\editionEtape{Le plus probable}Les secteurs sont équiprobables : « perdu », avec $\res6$ secteurs sur 10, domine.\editionEtape{La sucette}Trois secteurs sur dix, moins de la moitié : $\res{\text{non}}$, c’est peu probable, pas une chance sur deux.\editionEtape{Les ballons}Chaque taille accepte 5 couleurs : $2\times5=\res{10}$ ballons différents.
\end{blocEdition}
